\renewcommand{\thesection}{\arabic{section}}\renewcommand{\theHsection}{chapter.\arabic{section}}\setcounter{section}{15}\section{反常积分}
\noindent\begin{minipage}{\textwidth}\ProblemHeading{例题}{16.1（改写）}
设 $a>0$。证明 $\displaystyle\int_a^{+\infty}\frac{dx}{\sqrt[3]{x^2+x^4}}$ 收敛。
\end{minipage}\par\vfill
\noindent\begin{minipage}{\textwidth}\ProblemHeading{例题}{16.2（改写）}
证明 $\displaystyle\int_1^{+\infty}\frac{x^2}{x^4-x^2+1}\,dx$ 收敛。
\end{minipage}\par\vfill
\noindent\begin{minipage}{\textwidth}\ProblemHeading{例题}{16.3}
研究 $\displaystyle\int_1^{+\infty}\frac{\arctan x}{(1+x^2)^{3/2}}\,dx$ 的敛散性。
\end{minipage}\par\vfill
\newpage
\noindent\begin{minipage}{\textwidth}\ProblemHeading{例题}{16.4（改写）}
考察 $\displaystyle\int_0^{+\infty}\cos x\,dx$ 与 $\displaystyle\sum_{n=0}^{\infty}\int_{n\pi}^{(n+1)\pi}\cos x\,dx$ 的敛散性，说明将无穷积分分段后所得级数收敛，是否足以推出原积分收敛。
\end{minipage}\par\vfill
\noindent\begin{minipage}{\textwidth}\ProblemHeading{例题}{16.5（改写）}
设 $f(x)=x/(1+x^6\sin^2x)$（$x\ge0$）。
\SubQ{1}{证明 $\displaystyle\int_0^{+\infty}f(x)\,dx$ 收敛。}
\SubQ{2}{证明 $f$ 在 $[0,+\infty)$ 上无界。由此说明无穷积分收敛并不保证被积函数在无穷远处趋于 $0$。}
\end{minipage}\par\vfill
\noindent\begin{minipage}{\textwidth}\ProblemHeading{习题}{16.6}
判断下列无穷积分的敛散性（未注明范围的参数取实数）：
\SubQ{1}{$\displaystyle\int_1^{+\infty}\frac{x}{1+x^2}\,dx$。}
\SubQ{2}{$\displaystyle\int_0^{+\infty}\frac{3x^3-2}{x^5-x^3+1}\,dx$。}
\SubQ{3}{$\displaystyle\int_1^{+\infty}\frac{dx}{x\sqrt{x^2+1}}$。}
\end{minipage}\par\vfill
\newpage
\noindent\begin{minipage}{\textwidth}\ProblemHeading{习题}{16.7}
判断下列无穷积分的敛散性（未注明范围的参数取实数）：
\SubQ{1}{$\displaystyle\int_e^{+\infty}\frac{dx}{x(\ln x)^p}$。}
\SubQ{2}{$\displaystyle\int_1^{+\infty}x^{s-1}e^{-x}\,dx$。}
\SubQ{3}{$\displaystyle\int_1^{+\infty}\frac{(\ln x)^p}{1+x^2}\,dx$，$p>0$。}
\end{minipage}\par\vfill
\noindent\begin{minipage}{\textwidth}\ProblemHeading{习题}{16.8}
设 $\int_a^{+\infty}f(x)\,dx$ 收敛。证明：若 $\lim_{x\to+\infty}f(x)=b$ 存在，则 $b=0$。
\end{minipage}\par\vfill
\noindent\begin{minipage}{\textwidth}\ProblemHeading{习题}{16.9}
设 $\alpha>4$。证明 $\displaystyle\int_0^{+\infty}\frac{x}{1+x^\alpha\cos^2x}\,dx$ 收敛。
\end{minipage}\par\vfill
\newpage
\noindent\begin{minipage}{\textwidth}\ProblemHeading{习题}{16.10}
设 $f$ 在每个有限区间 $[a,A]$ 上可积。证明：$\int_a^{+\infty}f(x)\,dx$ 收敛于 $I$ 的充要条件是，对每个递增趋于 $+\infty$ 的数列 $\{A_n\}$（$A_1=a$），级数 $\displaystyle\sum_{n=1}^{\infty}\int_{A_n}^{A_{n+1}}f(x)\,dx$ 都收敛于 $I$。
\end{minipage}\par\vfill
\noindent\begin{minipage}{\textwidth}\ProblemHeading{例题}{16.11}
证明 $\displaystyle\int_1^{+\infty}\frac{\sin x}{x}\,dx$ 收敛，而 $\displaystyle\int_1^{+\infty}\frac{|\sin x|}{x}\,dx$ 发散。
\end{minipage}\par\vfill
\noindent\begin{minipage}{\textwidth}\ProblemHeading{例题}{16.12}
讨论 $\displaystyle\int_1^{+\infty}(-1)^{\lfloor x\rfloor}\,dx$ 和 $\displaystyle\int_1^{+\infty}(-1)^{\lfloor x^2\rfloor}\,dx$ 的敛散性，并在收敛时判断是否绝对收敛；$\lfloor t\rfloor$ 表示 $t$ 的整数部分。
\end{minipage}\par\vfill
\newpage
\noindent\begin{minipage}{\textwidth}\ProblemHeading{例题}{16.13}
设 $\max(p,q)>1$。证明 $\displaystyle\int_1^{+\infty}\frac{x\cos x}{x^p+x^q}\,dx$ 收敛。
\end{minipage}\par\vfill
\noindent\begin{minipage}{\textwidth}\ProblemHeading{习题}{16.14}
研究下列积分的敛散性，其中 $\lfloor t\rfloor$ 表示整数部分：
\SubQ{1}{$\displaystyle\int_0^{+\infty}\frac{\sqrt x\cos x}{1+x}\,dx$。}
\SubQ{2}{$\displaystyle\int_0^{+\infty}\frac{\sin((a+2)x)}{1+x^\alpha}\,dx$，$a\in\mathbb R$，$\alpha>0$。}
\SubQ{3}{$\displaystyle\int_1^{+\infty}\frac{\cos^2x}{x}\,dx$。}
\end{minipage}\par\vfill
\noindent\begin{minipage}{\textwidth}\ProblemHeading{习题}{16.15}
研究下列积分的敛散性，其中 $\lfloor t\rfloor$ 表示整数部分：
\SubQ{1}{$\displaystyle\int_0^{+\infty}\frac{\lfloor t\rfloor-t+1/2}{t+x}\,dt$，$x>0$。}
\SubQ{2}{$\displaystyle\int_0^{+\infty}\frac{\lfloor t\rfloor-t+a}{t+x}\,dt$，$x>0$，$a\ne1/2$。}
\end{minipage}\par\vfill
\newpage
\noindent\begin{minipage}{\textwidth}\ProblemHeading{习题}{16.16}
研究下列积分的绝对收敛性和条件收敛性：
\SubQ{1}{$\displaystyle\int_0^{+\infty}\frac{\sqrt x\sin x}{1+x}\,dx$。}
\SubQ{2}{$\displaystyle\int_1^{+\infty}\frac{\cos(1-2x)}{\sqrt{x^3}\sqrt[3]{x^2+1}}\,dx$。}
\SubQ{3}{$\displaystyle\int_0^{+\infty}\frac{\sin x}{\sqrt[3]{x^2+x+1}}\,dx$。}
\end{minipage}\par\vfill
\noindent\begin{minipage}{\textwidth}\ProblemHeading{习题}{16.17}
研究下列积分的绝对收敛性和条件收敛性：
\SubQ{1}{$\displaystyle\int_2^{+\infty}\frac{\sin x}{x\ln x}\,dx$。}
\end{minipage}\par\vfill
\noindent\begin{minipage}{\textwidth}\ProblemHeading{习题}{16.18}
设 $f$ 非负、单调递减，且 $\int_a^{+\infty}f(x)\,dx$ 收敛。证明 $f(x)=o(1/x)$（$x\to+\infty$）。
\end{minipage}\par\vfill
\newpage
\noindent\begin{minipage}{\textwidth}\ProblemHeading{习题}{16.19}
设 $f$ 在 $[a,+\infty)$ 上单调递减趋于 $0$。用 Dirichlet 判别法证明 $\int_a^{+\infty}f(x)\,dx$ 与 $\int_a^{+\infty}f(x)\sin^2x\,dx$ 同敛散。
\end{minipage}\par\vfill
\noindent\begin{minipage}{\textwidth}\ProblemHeading{习题}{16.20}
设 $f,g$ 在 $[a,b]$ 上连续，$g$ 在 $(a,b)$ 内可微且 $g^\prime(x)\ge0$。用分部积分证明：存在 $\xi\in[a,b]$，使\[\int_a^bf(x)g(x)\,dx=g(a)\int_a^\xi f(x)\,dx+g(b)\int_\xi^bf(x)\,dx.\]
\end{minipage}\par\vfill
\noindent\begin{minipage}{\textwidth}\ProblemHeading{习题}{16.21（校订）}
设 $P_m,P_n$ 分别为 $m,n$ 次非零多项式，且当 $x\ge a$ 时 $P_n(x)>0$。研究 $\displaystyle\int_a^{+\infty}\frac{P_m(x)}{P_n(x)}\sin x\,dx$ 的绝对收敛性和条件收敛性。
\end{minipage}\par\vfill
\newpage
\noindent\begin{minipage}{\textwidth}\ProblemHeading{习题}{16.22}
设 $f\in C^1([a,+\infty))$，且 $\int_a^{+\infty}f(x)\,dx$ 与 $\int_a^{+\infty}f^\prime(x)\,dx$ 都收敛。证明 $\lim_{x\to+\infty}f(x)=0$。
\end{minipage}\par\vfill
\noindent\begin{minipage}{\textwidth}\ProblemHeading{习题}{16.23}
设 $0<\alpha\le1/2$。证明 $\displaystyle\int_1^{+\infty}\frac{\sin x}{x^\alpha+\sin x}\,dx$ 发散。
\end{minipage}\par\vfill
\noindent\begin{minipage}{\textwidth}\ProblemHeading{习题}{16.24（校订）}
设 $f$ 在每个有限区间 $[0,b]$（$b>0$）上可积，且 $\lim_{x\to+\infty}f(x)=a$。证明\[\lim_{t\to0^+}t\int_0^{+\infty}e^{-tx}f(x)\,dx=a.\]
\end{minipage}\par\vfill
\newpage
\noindent\begin{minipage}{\textwidth}\ProblemHeading{例题}{16.25（改写）}
研究瑕积分 $\displaystyle\int_0^1\frac{dx}{\sqrt x}$ 的敛散性。
\end{minipage}\par\vfill
\noindent\begin{minipage}{\textwidth}\ProblemHeading{习题}{16.26（校订）}
设 $a<b$，$f,g$ 在每个 $[a+\varepsilon,b]$（$0<\varepsilon<b-a$）上可积。证明下列瑕积分判别结论：
\SubQ{1}{若靠近 $a$ 的右侧有 $0\le f(x)\le g(x)$，则 $\int_a^bg(x)\,dx$ 收敛蕴含 $\int_a^bf(x)\,dx$ 收敛；$\int_a^bf(x)\,dx$ 发散蕴含 $\int_a^bg(x)\,dx$ 发散。}
\SubQ{2}{若 $f\ge0$，$g$ 在 $a$ 的某个右邻域内为正，且 $\lim_{x\to a^+}f(x)/g(x)=l$，则 $0<l<+\infty$ 时两积分同敛散；$l=0$ 时 $g$ 的积分收敛蕴含 $f$ 的积分收敛；$l=+\infty$ 时 $g$ 的积分发散蕴含 $f$ 的积分发散。}
\SubQ{3}{$\int_a^bf(x)\,dx$ 收敛当且仅当对每个 $\epsilon>0$，存在 $0<\delta<b-a$，使 $0<\eta,\eta^\prime<\delta$ 时有 $\left|\int_{a+\eta}^{a+\eta^\prime}f(x)\,dx\right|<\epsilon$。}
\end{minipage}\par\vfill
\noindent\begin{minipage}{\textwidth}\ProblemHeading{习题}{16.27（校订）}
设 $a<b$，$f,g$ 在每个 $[a+\varepsilon,b]$（$0<\varepsilon<b-a$）上可积。证明下列瑕积分判别结论：
\SubQ{1}{若 $\int_a^b|f(x)|\,dx$ 收敛，则 $\int_a^bf(x)\,dx$ 收敛。}
\end{minipage}\par\vfill
\newpage
\noindent\begin{minipage}{\textwidth}\ProblemHeading{例题}{16.28}
研究 $\displaystyle\int_0^1\frac{\ln x}{1-x^2}\,dx$ 的敛散性。
\end{minipage}\par\vfill
\noindent\begin{minipage}{\textwidth}\ProblemHeading{例题}{16.29}
研究 $\displaystyle\int_0^1\frac{dx}{\sqrt[4]{1-x^4}}$ 的敛散性。
\end{minipage}\par\vfill
\noindent\begin{minipage}{\textwidth}\ProblemHeading{例题}{16.30}
研究 $\displaystyle\int_0^1x^{p-1}(1-x)^{q-1}\,dx$ 的敛散性，其中 $p,q\in\mathbb R$。
\end{minipage}\par\vfill
\newpage
\noindent\begin{minipage}{\textwidth}\ProblemHeading{例题}{16.31}
研究 $\displaystyle\int_0^{+\infty}x^{s-1}e^{-x}\,dx$ 的敛散性，其中 $s\in\mathbb R$。
\end{minipage}\par\vfill
\noindent\begin{minipage}{\textwidth}\ProblemHeading{例题}{16.32}
设 $p>0$。讨论 $\displaystyle\int_0^1\frac{\sin(1/x)}{x^p}\,dx$ 的敛散性。
\end{minipage}\par\vfill
\noindent\begin{minipage}{\textwidth}\ProblemHeading{例题}{16.33}
设 $\alpha>0$。讨论 $\displaystyle\int_0^{+\infty}\left[\left(1-\frac{\sin x}{x}\right)^{-\alpha}-1\right]dx$ 的敛散性。
\end{minipage}\par\vfill
\newpage
\noindent\begin{minipage}{\textwidth}\ProblemHeading{例题}{16.34（改写）}
考察 $\int_{-\infty}^{+\infty}x\,dx$ 的收敛性，并计算其 Cauchy 主值 $\displaystyle\lim_{A\to+\infty}\int_{-A}^Ax\,dx$。
\end{minipage}\par\vfill
\noindent\begin{minipage}{\textwidth}\ProblemHeading{例题}{16.35（改写）}
考察 $\displaystyle\int_{-\infty}^{+\infty}\frac{x}{1+x^2}\,dx$ 的收敛性，并计算其 Cauchy 主值 $\displaystyle\lim_{A\to+\infty}\int_{-A}^A\frac{x}{1+x^2}\,dx$。
\end{minipage}\par\vfill
\noindent\begin{minipage}{\textwidth}\ProblemHeading{例题}{16.36（改写）}
考察 $\displaystyle\int_{-1}^1\frac{dx}{x}$ 的收敛性，并计算其 Cauchy 主值 $\displaystyle\lim_{\varepsilon\to0^+}\left(\int_{-1}^{-\varepsilon}\frac{dx}{x}+\int_\varepsilon^1\frac{dx}{x}\right)$。
\end{minipage}\par\vfill
\newpage
\noindent\begin{minipage}{\textwidth}\ProblemHeading{习题}{16.37（校订）}
判断下列反常积分的敛散性（未注明范围的参数取实数）：
\SubQ{1}{$\displaystyle\int_0^{+\infty}\frac{x^\alpha}{1+x^\beta}\,dx$，$\beta\ge0$。}
\SubQ{2}{$\displaystyle\int_0^{+\infty}\frac{\ln(1+x)}{x^\alpha}\,dx$。}
\SubQ{3}{$\displaystyle\int_0^{+\infty}\frac{dx}{x^p+x^q}$。}
\end{minipage}\par\vfill
\noindent\begin{minipage}{\textwidth}\ProblemHeading{习题}{16.38（校订）}
判断下列反常积分的敛散性（未注明范围的参数取实数）：
\SubQ{1}{$\displaystyle\int_0^1\frac{\sqrt x}{e^{\sin x}-1}\,dx$。}
\SubQ{2}{$\displaystyle\int_0^{\pi/2}\frac{\ln(\sin x)}{\sqrt x}\,dx$。}
\SubQ{3}{$\displaystyle\int_0^{+\infty}\frac{dx}{x^p|\ln x|^q}$。}
\end{minipage}\par\vfill
\noindent\begin{minipage}{\textwidth}\ProblemHeading{习题}{16.39（校订）}
判断下列反常积分的敛散性（未注明范围的参数取实数）：
\SubQ{1}{$\displaystyle\int_0^{+\infty}\sin(x^\alpha)\,dx$，$\alpha>1$。}
\end{minipage}\par\vfill
\newpage
\noindent\begin{minipage}{\textwidth}\ProblemHeading{习题}{16.40}
判断下列反常积分的绝对收敛性和条件收敛性：
\SubQ{1}{$\displaystyle\int_0^{+\infty}x^p\sin(x^q)\,dx$，$p\in\mathbb R$，$q\ne0$。}
\SubQ{2}{$\displaystyle\int_0^{+\infty}\frac{x^p\sin x}{1+x^q}\,dx$，$p\in\mathbb R$，$q\ge0$。}
\end{minipage}\par\vfill
\noindent\begin{minipage}{\textwidth}\ProblemHeading{习题}{16.41}
判断 $\displaystyle\int_0^{+\infty}\frac{\sin(x+1/x)}{x^\alpha}\,dx$ 的绝对收敛性和条件收敛性，其中 $\alpha\in\mathbb R$。
\end{minipage}\par\vfill
\noindent\begin{minipage}{\textwidth}\ProblemHeading{习题}{16.42}
设 $f$ 在 $(0,1)$ 上单调，在两个端点附近未必有界，且 $\int_0^1f(x)\,dx$ 收敛。证明\[\lim_{n\to\infty}\frac1n\sum_{k=1}^{n-1}f(k/n)=\int_0^1f(x)\,dx.\]
\end{minipage}\par\vfill
\newpage
\noindent\begin{minipage}{\textwidth}\ProblemHeading{习题}{16.43}
设 $f$ 在 $(0,1)$ 上单调，且 $\displaystyle\lim_{n\to\infty}\frac1n\sum_{k=1}^{n-1}f(k/n)$ 存在，是否能保证 $\int_0^1f(x)\,dx$ 收敛？试研究 $f(x)=\dfrac{1}{x}-\dfrac{1}{1-x}$。
\end{minipage}\par\vfill
\noindent\begin{minipage}{\textwidth}\ProblemHeading{习题}{16.44}
利用单调函数的等距 Riemann 和收敛到其收敛反常积分这一结论，证明：对任意 $\alpha>0$，有\[\lim_{n\to\infty}\frac{1^{\alpha-1}+2^{\alpha-1}+\cdots+n^{\alpha-1}}{n^\alpha}=\frac1\alpha.\]
\end{minipage}\par\vfill
\noindent\begin{minipage}{\textwidth}\ProblemHeading{习题}{16.45}
设 $f$ 在 $\mathbb R$ 上连续，且 $\lim_{x\to+\infty}f(x)=a$、$\lim_{x\to-\infty}f(x)=b$ 均有限。证明对任意 $\eta>0$，\[\int_{-\infty}^{+\infty}[f(x+\eta)-f(x)]\,dx=\eta(a-b).\]
\end{minipage}\par\vfill
\newpage
\noindent\begin{minipage}{\textwidth}\ProblemHeading{习题}{16.46（校订）}
设 $a,b>0$，证明下列公式，其中 $f(+\infty)$ 表示 $f$ 在 $+\infty$ 处的有限极限：
\SubQ{1}{若 $f$ 在 $[0,+\infty)$ 上连续且 $f(+\infty)$ 存在，则 $\displaystyle\int_0^{+\infty}\frac{f(ax)-f(bx)}x\,dx=[f(0)-f(+\infty)]\ln\frac ba$。}
\SubQ{2}{若 $f$ 在 $[0,+\infty)$ 上连续且 $\int_1^{+\infty}f(x)/x\,dx$ 收敛，则 $\displaystyle\int_0^{+\infty}\frac{f(ax)-f(bx)}x\,dx=f(0)\ln\frac ba$。}
\SubQ{3}{若 $f$ 在 $(0,+\infty)$ 上连续，$f(+\infty)$ 存在，且 $\int_0^1f(x)/x\,dx$ 收敛，则 $\displaystyle\int_0^{+\infty}\frac{f(ax)-f(bx)}x\,dx=f(+\infty)\ln\frac ab$。}
\end{minipage}\par\vfill
\noindent\begin{minipage}{\textwidth}\ProblemHeading{习题}{16.47}
设 $a,b>0$。利用以下 Frullani 公式计算下列积分：若 $f$ 在 $[0,+\infty)$ 上连续且 $f(+\infty)=l$ 存在，则\[\int_0^{+\infty}\frac{f(ax)-f(bx)}x\,dx=[f(0)-l]\ln\frac ba;\]若改为假设 $\int_1^{+\infty}f(x)/x\,dx$ 收敛，则右端为 $f(0)\ln(b/a)$。
\SubQ{1}{$\displaystyle\int_0^{+\infty}\frac{e^{-ax}-e^{-bx}}x\,dx$。}
\SubQ{2}{$\displaystyle\int_0^{+\infty}\frac{\cos(ax)-\cos(bx)}x\,dx$。}
\end{minipage}\par\vfill
\noindent\begin{minipage}{\textwidth}\ProblemHeading{例题}{16.48}
设 $a>0$，$D=\{(x,y):x^2+y^2\ge a^2\}$。讨论 $\displaystyle\iint_D\frac{dx\,dy}{(x^2+y^2)^{p/2}}$ 的敛散性，其中 $p\in\mathbb R$。
\end{minipage}\par\vfill
\newpage
\noindent\begin{minipage}{\textwidth}\ProblemHeading{例题}{16.49（校订）}
设 $f(x,y)=(2-xy)xye^{-xy}$，$D=\{(x,y):0\le x\le1,\ y\ge0\}$，$D_R=[0,1]\times[0,R]$。
\SubQ{1}{计算 $\displaystyle\lim_{R\to+\infty}\iint_{D_R}f(x,y)\,dx\,dy$。}
\SubQ{2}{分别计算 $\displaystyle\int_0^1dx\int_0^{+\infty}f(x,y)\,dy$ 与 $\displaystyle\int_0^{+\infty}dy\int_0^1f(x,y)\,dx$。}
\SubQ{3}{判断 $\iint_D|f(x,y)|\,dx\,dy$ 是否收敛，并据此判断按任意截断区域定义的反常重积分 $\iint_Df(x,y)\,dx\,dy$ 是否收敛。}
\end{minipage}\par\vfill
\noindent\begin{minipage}{\textwidth}\ProblemHeading{例题}{16.50}
设 $a>0$、$p>0$，$D=\{(x,y):x^2+y^2\le a^2\}$。讨论 $\displaystyle\iint_D\frac{dx\,dy}{(x^2+y^2)^{p/2}}$ 的敛散性。
\end{minipage}\par\vfill
\noindent\begin{minipage}{\textwidth}\ProblemHeading{例题}{16.51}
计算 $\displaystyle I=\iint_D\ln\sin(x-y)\,dx\,dy$，其中 $D$ 是由 $y=0$、$y=x$ 与 $x=\pi$ 围成的区域。
\end{minipage}\par\vfill
\newpage
\noindent\begin{minipage}{\textwidth}\ProblemHeading{例题}{16.52}
设 $D=\{(x,y,z):x>0,y>0,z>0,x+y+z<1\}$。研究 $\displaystyle\iiint_D\frac{dx\,dy\,dz}{(x+y+z)^p}$ 的敛散性，其中 $p\in\mathbb R$。
\end{minipage}\par\vfill
\noindent\begin{minipage}{\textwidth}\ProblemHeading{习题}{16.53}
讨论下列反常重积分的敛散性（未注明范围的参数取实数）：
\SubQ{1}{$\displaystyle\iint_{\mathbb R^2}\frac{dx\,dy}{(1+|x|^p)(1+|y|^q)}$。}
\SubQ{2}{$\displaystyle\iint_{0<y<1}\frac{dx\,dy}{(1+x^2+y^2)^p}$，其中 $x\in\mathbb R$。}
\SubQ{3}{$\displaystyle\iint_D\frac{dx\,dy}{(x^\alpha+y^\beta)^m}$，$\alpha,\beta>0$，$D=\{x\ge0,y\ge0,x^\alpha+y^\beta\ge1\}$。}
\end{minipage}\par\vfill
\noindent\begin{minipage}{\textwidth}\ProblemHeading{习题}{16.54}
讨论下列反常重积分的敛散性（未注明范围的参数取实数）：
\SubQ{1}{$\displaystyle\iint_D\frac{dx\,dy}{(x+y)^p}$，$D=\{0\le x\le1,x+y\ge1\}$。}
\SubQ{2}{$\displaystyle\iint_{x^2+y^2<1}\frac{dx\,dy}{(1-x^2-y^2)^p}$。}
\SubQ{3}{$\displaystyle\iint_D\frac{dx\,dy}{(1-x^\alpha-y^\beta)^p}$，$\alpha,\beta>0$，$D=\{x\ge0,y\ge0,x^\alpha+y^\beta<1\}$。}
\end{minipage}\par\vfill
\newpage
\noindent\begin{minipage}{\textwidth}\ProblemHeading{习题}{16.55}
讨论下列反常重积分的敛散性（未注明范围的参数取实数）：
\SubQ{1}{$\displaystyle\iint_D\frac{x-y}{(x+y)^3}\,dx\,dy$，$D=[0,1]^2$。}
\SubQ{2}{$\displaystyle\iint_D\frac{(x+y)^2\ln(x+y)}{x^2+y^2}\,dx\,dy$，$D=\{x\ge0,y\ge0,x+y\le1\}$。}
\end{minipage}\par\vfill
\noindent\begin{minipage}{\textwidth}\ProblemHeading{习题}{16.56}
证明下列结论：
\SubQ{1}{$\displaystyle\int_0^{+\infty}dy\int_0^{+\infty}\sin(x^2+y^2)\,dx=\int_0^{+\infty}dx\int_0^{+\infty}\sin(x^2+y^2)\,dy=\frac\pi4$。}
\SubQ{2}{反常重积分 $\displaystyle\iint_D\sin(x^2+y^2)\,dx\,dy$ 发散，其中 $D=\{(x,y):x\ge0,y\ge0\}$。}
\end{minipage}\par\vfill
\noindent\begin{minipage}{\textwidth}\ProblemHeading{习题}{16.57（校订）}
计算下列反常重积分；含参数的小问先判断其收敛范围，再在收敛时求值：
\SubQ{1}{$\displaystyle\iint_D\frac{dx\,dy}{(1+x+y)^3}$，$D=\{x\ge0,y\ge0\}$。}
\SubQ{2}{$\displaystyle\iint_D\frac{dx\,dy}{e^x+e^y}$，$D=\{x\ge0,y\ge0\}$。}
\SubQ{3}{$\displaystyle\iint_D\frac{dx\,dy}{(x+y)^p}$，$p>1$，$D=\{0\le x\le1,x+y\ge1\}$。}
\end{minipage}\par\vfill
\newpage
\noindent\begin{minipage}{\textwidth}\ProblemHeading{习题}{16.58（校订）}
计算下列反常重积分；含参数的小问先判断其收敛范围，再在收敛时求值：
\SubQ{1}{$\displaystyle\iint_D\frac{dx\,dy}{x^py^q}$，$p>q>1$，$D=\{x\ge1,xy\ge1\}$。}
\SubQ{2}{$\displaystyle\iint_{x^2+y^2<1}\frac{dx\,dy}{\sqrt{1-x^2-y^2}}$。}
\SubQ{3}{$\displaystyle\iint_De^{-(x+y)}\,dx\,dy$，$D=\{0\le x\le y\}$。}
\end{minipage}\par\vfill
\noindent\begin{minipage}{\textwidth}\ProblemHeading{习题}{16.59（校订）}
计算下列反常重积分；含参数的小问先判断其收敛范围，再在收敛时求值：
\SubQ{1}{$\displaystyle\iint_{\mathbb R^2}e^{-(x^2+y^2)}\cos(x^2+y^2)\,dx\,dy$。}
\SubQ{2}{$\displaystyle\iint_{\mathbb R^2}e^{-(x^2+y^2)}\sin(x^2+y^2)\,dx\,dy$。}
\SubQ{3}{$\displaystyle\iint_{x^2+y^2\le1}\ln\frac1{\sqrt{x^2+y^2}}\,dx\,dy$。}
\end{minipage}\par\vfill
\noindent\begin{minipage}{\textwidth}\ProblemHeading{习题}{16.60（校订）}
计算下列反常重积分；含参数的小问先判断其收敛范围，再在收敛时求值：
\SubQ{1}{$\displaystyle\iint_D\ln\sin(x-y)\,dx\,dy$，$D=\{0\le y\le x\le\pi\}$。}
\SubQ{2}{$\displaystyle\iiint_D\frac{dx\,dy\,dz}{(x+y+z)^p}$，$p>1$，$D=\{x\ge0,y\ge0,z\ge0,x+y+z\le1\}$。}
\end{minipage}\par\vfill
